This work investigates Breadth-First Search execution on GPUs using CUDA and
NVSHMEM, with Graph500 as the reference for graph generation, graph construction,
and performance evaluation. Two implementations were developed: a CUDA/NVSHMEM
version compatible with the Graph500 workflow and a refactored graph traversal
library executed mostly on the GPU. The evaluation used a personal computer with
one NVIDIA RTX 5060 Ti GPU and a GB200 environment with four GPUs. The results
show significant improvements in R-MAT generation and BFS performance at larger
scales. On the personal computer, GPU graph generation was approximately $7.4$
times faster than the CPU reference at \texttt{SCALE=22}, and the library reached
2.138 GTEPS at \texttt{SCALE=21}. On the GB200 environment, at
\texttt{SCALE=29}, the library reached 6.180 GTEPS, compared with 1.570 GTEPS for
the CPU reference, while reducing graph generation time from 31.13 s to 2.738 s.
The results indicate that keeping BFS coordination on the device is a promising
direction for distributed graph traversal, although irregularity, remote
communication, and memory usage remain limiting factors.

\textbf{Keywords}: Breadth-First Search, GPU, CUDA, NVSHMEM, Graph500.
